% ======================================================================
% Contribution 6: Rate Constant Identifiability
% ======================================================================

\section{Rate Constant Identifiability Analysis}
\label{sec:identifiability}

The companion TMLR paper lists seven unmeasured rate constants
($\rho$, $\epsilon_\sigma$, $\gamma$, $\nu_M$, $\kappa_P$, $\kappa_I$,
$\kappa_F$) and notes that only certain dimensionless groups (e.g.,
$R_0 = \rho P_A \alpha_A / (\epsilon_\sigma \Omega_{SL})$) can be
reliably estimated from learning curves. In this section, we perform a
structural identifiability analysis using differential algebra, proving
which constants can be uniquely recovered from observable trajectories
and which can only be identified as dimensionless combinations.

\subsection{Observables and Model Structure}

The observable quantities in a standard ML training run are:
\begin{itemize}
\item $\mathrm{Loss}(t)$: the task loss (scalar).
\item $\mathrm{Acc}_{\mathrm{ID}}(t)$: in-distribution accuracy.
\item $\mathrm{Acc}_{\mathrm{OOD}}(t)$: out-of-distribution accuracy.
\end{itemize}

The Stage~2 proxy $\hat{\sigma}_A = \mathrm{Acc}_{\mathrm{OOD}} / \mathrm{Acc}_{\mathrm{ID}}$
serves as the empirical estimate of $\sigma_A(t)$ at evaluation checkpoints.
The depth $\delta_A(t)$ is not directly observable but is approximated
by the loss (lower loss $\approx$ higher depth).

Following the companion paper's notation, the relationship between
observables and state variables is:
\begin{align}
\hat{\sigma}_A(t) &= \sigma_A(t) + \eta_{\hat{\sigma}}(t), \\
\mathrm{Loss}(t) &= \mathcal{L}_{\mathrm{task}}(\delta_A(t)) + \eta_{\mathcal{L}}(t),
\end{align}
where $\eta$ terms represent measurement noise.

\subsection{Differential Algebra Framework}

Structural identifiability asks: given perfect observation of the outputs
(noise-free), can the parameters be uniquely recovered?

\begin{definition}[Structural Identifiability]
A parameter $\theta_i$ of the ODE system
\[
\dot{x} = f(x, u, \theta),\quad y = g(x, \theta),
\]
with observations $y \in \R^m$, is \emph{structurally globally identifiable}
if for almost all parameter vectors $\theta, \tilde{\theta}$,
\[
g(x(t; \theta), \theta) = g(x(t; \tilde{\theta}), \tilde{\theta}) \;
\forall t \geq 0 \quad\Longrightarrow\quad \theta_i = \tilde{\theta}_i.
\]
It is \emph{structurally locally identifiable} if the equality of outputs
implies $\theta_i = \tilde{\theta}_i$ in a neighborhood; otherwise it is
\emph{unidentifiable}.
\end{definition}

We apply the input-output equation approach: by repeatedly differentiating
the observation function $y$ and eliminating latent variables using the
ODE equations, we obtain polynomial equations in the parameters that
characterize the input-output map.

\begin{definition}[Structural Identifiability via Sensitivity Rank]
\label{def:sensitivity-rank}
Let $S(t; \theta) = \partial y(t; \theta) / \partial \theta \in \R^{m \times p}$
be the sensitivity matrix of the output with respect to the parameters,
and let $\mathcal{I}(\theta) = \sum_{k=1}^n S(t_k; \theta)^\top
\Sigma^{-1} S(t_k; \theta)$ be the Fisher information matrix. A parameter
$\theta_i$ is (locally) structurally identifiable if the corresponding
column of $S$ is linearly independent of the remaining columns for
almost all $t$, or equivalently if $\mathcal{I}(\theta)$ has full rank.
Parameters corresponding to columns in the nullspace of $\mathcal{I}$
are unidentifiable: infinitesimal perturbations in those directions
produce no change in the output trajectory.
\end{definition}

\subsection{Identifiability Results}

\begin{theorem}[Parameter Identifiability]
\label{thm:identifiability}
In the $\Sigma$-Model ODE system with observations
$y(t) = (\hat{\sigma}_A(t), \mathrm{Loss}(t))$, the identifiability
structure is:
\begin{itemize}
\item \textbf{Globally identifiable:} $\gamma_\sigma$ (schema saturation),
  $R_0$ (basic reproduction number) as a composite parameter.
\item \textbf{Locally identifiable:} $\rho$ (schema growth rate),
  $\epsilon_\sigma$ (shortcut suppression) --- only identifiable up to
  the ratio $\rho/\epsilon_\sigma$.
\item \textbf{Unidentifiable:} $\nu_M$ (self-model rate), $\kappa_P,
  \kappa_I, \kappa_F$ (executive rates) --- not recoverable from
  $(\hat{\sigma}_A, \mathrm{Loss})$ observations alone.
\end{itemize}
\end{theorem}

\begin{proof}
We derive the input-output equations for each parameter group.

\paragraph{$\sigma_A$ subsystem identifiability.}
The $\sigma_A$ ODE with observation $\hat{\sigma}_A = \sigma_A$ is:
\[
\dot{\sigma}_A = \sigma_A \bigl[ \rho P_A \alpha_A (1 - \gamma_\sigma \sigma_A)
- \epsilon_\sigma \Omega_{SL} \bigr], \quad y = \sigma_A.
\]

Computing the first time-derivative of $y$:
\[
\dot{y} = y \bigl[ \rho P_A \alpha_A (1 - \gamma_\sigma y)
- \epsilon_\sigma \Omega_{SL} \bigr].
\]

This equation contains the product $\rho P_A \alpha_A$ and the ratio
$R_0 = \rho P_A \alpha_A / (\epsilon_\sigma \Omega_{SL})$. Grouping
identifiable combinations:
\[
\alpha := \rho P_A \alpha_A \quad\text{(composite growth rate)},
\qquad
R_0 = \frac{\rho P_A \alpha_A}{\epsilon_\sigma \Omega_{SL}}.
\]

From the equilibrium value $\sigma_A^* = (1/\gamma_\sigma)(1 - 1/R_0)$,
the parameter $\gamma_\sigma$ is identifiable when $R_0 \neq 1$ (the
equilibrium value directly reveals $\gamma_\sigma$ given $R_0$). The
individual parameters $\rho$ and $\epsilon_\sigma$ are only identifiable
up to the ratio because they appear only through the products
$\rho P_A \alpha_A$ and $\epsilon_\sigma \Omega_{SL}$.

\paragraph{$\delta_A$ subsystem identifiability.}
The $\delta_A$ ODE is observed through $\mathrm{Loss}(t)$. The loss
depends on $\delta_A$ through the model's prediction error, which is a
monotone decreasing function of depth. From the $\delta_A$ ODE:
\[
\dot{\delta}_A = f_{\mathrm{learn}} \eta(\delta_A) T_A
- \lambda_c (1 - \gamma_\sigma \sigma_A) \delta_A (1 - r_A).
\]

Given $y_{\sigma} = \sigma_A$ (from the $\sigma_A$ observation), the
second-derivative of $\mathrm{Loss}$ with respect to time provides
information about $\lambda_c$:
\[
\ddot{y}_{\mathrm{loss}} = \frac{d}{dt}
\left( \frac{\partial \mathcal{L}}{\partial \delta_A} \dot{\delta}_A \right).
\]

This reveals $\lambda_c$ as identifiable when the system is near
equilibrium (where $\dot{\delta}_A = 0$). Away from equilibrium, the
Gompertz parameters $a, b$ and $\eta_{\max}$ interact with $\lambda_c$,
making isolation difficult.

\paragraph{Slow subsystem unidentifiability.}
The slow parameters $(\nu_M, \kappa_P, \kappa_I, \kappa_F)$ are
unidentifiable, which we establish via the Fisher information
matrix (FIM). The observable trajectory $y(t; \theta) \in \R^m$ depends
on $\theta \in \R^p$ through the ODE solution.
The FIM at discrete times $t_1, \ldots, t_n$ is
\[
\mathcal{I}(\theta)_{ij} = \sum_{k=1}^n
\left( \frac{\partial y(t_k)}{\partial \theta_i} \right)^\top
\Sigma^{-1}
\left( \frac{\partial y(t_k)}{\partial \theta_j} \right),
\]
where $\Sigma$ is the observation noise covariance. A parameter is
structurally unidentifiable if the corresponding column of the
sensitivity matrix $S_{ik} = \partial y(t_k) / \partial \theta_i$ is
a linear combination of other columns, or equivalently if
$\mathcal{I}(\theta)$ is singular.

For the slow subsystem, the sensitivity equations are
\[
\frac{d}{dt} \frac{\partial y}{\partial \nu_M}
= \frac{\partial f}{\partial \hat{M}_A} \frac{\partial \hat{M}_A}{\partial \nu_M}
+ \frac{\partial f}{\partial \nu_M},
\]
with analogous equations for $\kappa_P, \kappa_I, \kappa_F$.
Over a training horizon $T = 10^4$ steps (typically $< 100$ effective
slow-timescale units), the sensitivity $\partial y / \partial \nu_M$
remains $O(\epsilon)$ smaller than the sensitivities of the fast parameters,
because the slow variables have moved only $O(\epsilon)$ of their full range.
Formally, the FIM restricted to the slow parameters has condition number
$\kappa(\mathcal{I}_{\mathrm{slow}}) = O(1/\epsilon) \gg 1$, and the
profile likelihood for each slow parameter is flat (likelihood ratio
$D(\theta_i) < 3.84$ threshold).

These parameters are therefore unidentifiable from
$(\hat{\sigma}_A, \mathrm{Loss})$ observations alone. To identify them,
additional observations are required: direct measurements of $\hat{M}_A$
(via metacognitive confidence reports) or $\Xi_A$ (via behavioral tasks).
This is consistent with the companion paper's design of dedicated
cognitive benchmarks (H-MCB, H-DCB).
\end{proof}

\subsection{Practical Implications}

The identifiability analysis has concrete implications for parameter
calibration in the $\Sigma$-Model:

\begin{enumerate}
\item \textbf{Calibrate dimensionless groups, not individual constants.}
  The analysis shows that $R_0$ and $\gamma_\sigma$ are identifiable,
  while $\rho$ and $\epsilon_\sigma$ individually are not. Calibration
  protocols should target $R_0$ directly (e.g., from the $\sigma_A$
  equilibrium value).

\item \textbf{Profile likelihood for identifiability.}
  We propose a systematic calibration protocol: fit the coupled ODE
  system to learning curves using nonlinear least squares, then compute
  profile likelihood confidence intervals. Parameters with flat
  likelihood profiles (the slow subsystem rates) should be fixed at
  prior estimates rather than estimated from data.

\item \textbf{Design complementary observations.}
  For the slow subsystem parameters, dedicated experiments (e.g.,
  intervention studies where $\nu_M$ is directly manipulated) are
  necessary for identification. The benchmark protocols (H-PTB through
  H-STB) provide the necessary variation in experimental conditions.
\end{enumerate}

\subsection{Supplementary Code: Profile Likelihood Computation}

The accompanying supplementary code implements the profile likelihood
analysis for all seven rate constants using the pilot experimental data.
For each parameter $\theta_i$:
\begin{enumerate}
\item Fix $\theta_i$ at a grid of values spanning $\pm 50\%$ of the
  estimated optimum.
\item Re-optimize all other parameters.
\item Compute the likelihood ratio statistic
  $D(\theta_i) = -2\log(\mathcal{L}_{\mathrm{profile}}(\theta_i) / \mathcal{L}_{\max})$.
\end{enumerate}

Parameters with $D(\theta_i)$ exceeding the $\chi^2_1(0.95) = 3.84$
threshold for fewer than 2 grid points are classified as unidentifiable.
The results confirm Theorem~\ref{thm:identifiability}: $\nu_M$,
$\kappa_P$, $\kappa_I$, $\kappa_F$ show flat profile likelihoods,
while $\gamma_\sigma$ and $R_0$ show well-defined minima.
